\textbf{1. Kiến thức cần nhớ}
\begin{itemize}
\item \textbf{Phương trình tích}: Phương trình có dạng $A(x) \cdot B(x) = 0$.
$$\boxed{A(x) \cdot B(x) = 0 \iff \left[\begin{aligned} &A(x) = 0 \\ &B(x) = 0 \end{aligned}\right.}$$
\item \textbf{Phương trình chứa ẩn ở mẫu}: Là phương trình có chứa biến số ở mẫu thức.
\begin{itemize}
\item \textbf{Điều kiện xác định (ĐKXĐ)}: Là điều kiện để tất cả các mẫu thức trong phương trình đều khác $0$.
\item Phép quy đồng khử mẫu có thể làm xuất hiện nghiệm ngoại lai nên \textbf{bắt buộc phải đối chiếu ĐKXĐ} trước khi kết luận nghiệm.
\end{itemize}
\end{itemize}

\textbf{2. Phương pháp giải các dạng}

%% DANG 9D2?4-1
\textbf{Dạng 1. Phương trình tích}
\begin{enumerate}
\item Chuyển tất cả các hạng tử sang vế trái để vế phải bằng $0$.
\item Phân tích đa thức ở vế trái thành nhân tử (đặt nhân tử chung, dùng hằng đẳng thức hoặc tách hạng tử).
\item Cho từng nhân tử bằng $0$ để giải các phương trình bậc nhất thành phần.
\item Kết luận tập nghiệm $S$ (lấy tất cả các nghiệm tìm được).
\end{enumerate}

%% DANG 9D2?4-2
\textbf{Dạng 2. Phương trình chứa ẩn ở mẫu}
\begin{enumerate}
\item \textbf{Bước 1}: Tìm điều kiện xác định của phương trình (cho tất cả các mẫu khác $0$).
\item \textbf{Bước 2}: Phân tích các mẫu thành nhân tử để tìm Mẫu thức chung (MTC).
\item \textbf{Bước 3}: Quy đồng mẫu thức hai vế rồi khử mẫu.
\item \textbf{Bước 4}: Giải phương trình thu được.
\item \textbf{Bước 5}: \textbf{Đối chiếu với ĐKXĐ} để loại nghiệm không thỏa mãn và kết luận nghiệm.
\end{enumerate}

%% DANG 9D2?4-3
\textbf{Dạng 3. Giải bài toán bằng cách lập phương trình quy về bậc nhất}
\begin{enumerate}
\item Chọn một đại lượng làm ẩn số $x$ kèm đơn vị và điều kiện thích hợp.
\item Biểu diễn các đại lượng khác theo $x$ dưới dạng phân thức hoặc tích.
\item Lập phương trình và giải phương trình (quy đồng khử mẫu hoặc đưa về phương trình tích).
\item So sánh với điều kiện ban đầu và trả lời bài toán.
\end{enumerate}
